\documentclass[11pt]{article}
\usepackage[margin=1.5cm]{geometry}
\usepackage{amsmath,amssymb}
\usepackage{enumitem}

\begin{document}
\twocolumn

\title{Chapter 14 In-Class Exercise: Reading Student Data from a File}
\author{Prof. Jordan C. Hanson}
\date{\today}
\maketitle

\section{A Simple Course Class}

Suppose we want a \texttt{Course} object to store a course name and a list of students.

\begin{verbatim}
class Course:
    def __init__(self, name):
        self.name = name
        self.students = []
    def add_student(self,id,name):
        self.students.append([id, name])
\end{verbatim}

Create the course:

\begin{verbatim}
course = Course("COSC 120")
\end{verbatim}

At this point, what is stored in \texttt{course.students}?

\section{Student Data in a File}

Create a text file named

\begin{verbatim}
students.txt
\end{verbatim}

containing

\begin{verbatim}
2601001,Ana
2601002,Ben
2601003,Chloe
2601004,Diego
2601005,Elena
\end{verbatim}

\section{Open and Read the File}

Python can open a file with

\begin{verbatim}
with open("students.txt", "r") as f:
\end{verbatim}

The \texttt{"r"} means \textbf{read mode}.

Complete the loop:

\begin{verbatim}
with open("students.txt", "r") as f:
    for line in f:
        id, name = line.strip().split(",")
        course.add_student(int(id),name)
\end{verbatim}

\vspace{0.3cm}

\section{Examine the Course Object}

After the file has been read, run

\begin{verbatim}
print(course.name)
print(course.students)
\end{verbatim}

How many student records should now be stored?  Now print them one at a time:

\begin{verbatim}
for student in course.students:
    print(student)
\end{verbatim}

\section{Add a Useful Method}

Add this method to \texttt{Course}:

\begin{verbatim}
def number_of_students(self):
    return len(self.students)
\end{verbatim}

Then run

\begin{verbatim}
print(course.number_of_students())
\end{verbatim}

What should be printed?

\section{Your Turn}

Add another line to \texttt{students.txt}:

\begin{verbatim}
2601006,Fatima
\end{verbatim}

Run the program again.

\begin{enumerate}
\item Does the Python code need to be changed?
\item How many records are now stored?
\item Why is reading data from a file more useful than typing every
student directly into the Python program?
\end{enumerate}

\vspace{0.25cm}
\noindent
\textbf{Key idea:} The file stores the \emph{data}; the
\texttt{Course} class stores and organizes that data after Python
reads it.

\end{document}
